% ANCHOR_GAP_PAPER.tex — C41 Track-B Shape-2, the anchor-gap census paper.
% Lane: c41-trackb-shape2-mzr-20260804T1452K. Drafted by Lana 2026-08-04 (KST).
% Every number in this draft traces to a receipted lane artifact; the source →
% section map is PAPER_CHANGELOG.md. Compile requires AASTeX 6.3.1 (aastex631.cls);
% compilation was not run inside the lane (no-network / lane-only-writes rule).
\documentclass[twocolumn]{aastex702}

\newcommand{\te}{$T_e$}
\newcommand{\logm}{\log M_*/M_\odot}
\newcommand{\oh}{12+\log(\mathrm{O/H})}

\begin{document}

\title{The Public-Archive Direct-\te{} Anchor Gap at $z>3$: A Contract-Grade Census}

\author{NebulaMind Autonomous Research Crew}
\email{lab@nebulamind.net}
\affiliation{NebulaMind Laboratory (autonomous research system); drafting agent: Lana;
pipeline author: Hwao; adversarial review: Kun; receipts: Tori.
Lane \texttt{c41-trackb-shape2-mzr-20260804T1452K}.}

\begin{abstract}
Whether locally calibrated metallicity diagnostics survive at high redshift is a
disputed axis of the galaxy-evolution literature, and its named settle-line is
statistical: grow the $z>3$ auroral-line/direct-\te{} anchor set beyond the
$\sim$25 galaxies the corpus records, then re-test the disputed calibration claims on
matched-mass samples. We report a census of what the \emph{public} archives can
actually supply toward that settle-line, executed under a measurement contract frozen
before any science row was fetched: a declared single \te{}-anchored scale, a uniform
pre-registered S/N$(\lambda4363)\geq5$ floor, a frozen direct-method pipeline applied
to source-published fluxes, per-row exclusion provenance, and a pre-committed
forecast defining the null's information content. A global VizieR TAP enumeration
(79 tables carrying a $\lambda$4363-class column; 23 joinable to a redshift; 8
catalogs reachable at run time) yielded 95 rows at $z>3$ with a tabulated auroral
flux, of which exactly \textbf{5} survive as contract-grade direct-\te{} anchors with
linked stellar masses ($z=4.015$--$8.496$; $\oh = 7.109$--$8.032$). The frozen
forecast conservatively expected $\sim$25 usable anchors (its licensed pre-fetch
revision expected 87); the realized public yield falls short of even the conservative
expectation by roughly an order of magnitude. No stellar-mass bin reaches the
pre-committed 3-anchor minimum, so \emph{no deficit verdict of any size or direction
is possible at contract-grade public statistics}: the calibration-validity dispute
currently rests on data that are not publicly quotable at uniform rigor. This paper
is a census and a null. We tabulate the five anchors with full provenance, the
complete 90-row exclusion taxonomy (including a five-object near-miss band at
S/N$=4.76$--$4.84$, within 5\% of the floor), and the specific publication acts —
machine-readable flux tables with uncertainties, linked stellar masses, and
per-object magnifications — that would close the gap.
\end{abstract}

\keywords{galaxies: abundances --- galaxies: high-redshift --- methods: data analysis ---
astronomical databases: catalogs}

\section{Introduction}\label{sec:intro}

The debate this census serves is ledger-mapped, not paraphrased. In the crew's C41
status/debate map of the JWST high-$z$ galaxy-evolution literature, axis A3 —
``do locally calibrated metallicity diagnostics survive at high redshift?'' — is the
corpus's largest axis (22 entries) and its clearest cross-paper contradiction
cluster. One stance-verified side holds that strong-line calibrations transfer to
high redshift without strong evolution (ledger entry c41\_043); the opposing cluster
(c41\_033, c41\_035, c41\_044, c41\_045, c41\_049) holds that local calibrations
shift, disagree mutually, or fail off the local excitation locus, with reported
discrepancies from 0.05--0.1 dex (local-calibration underestimates) to a
0.35$\pm$0.28 dex scatter between UV- and optical-derived values. The axis's binding
capability statistic is ledger entry c41\_012: auroral-line detections at $z>3$
remain limited to $\sim$25 galaxies. The map's settle-line is explicit: grow the
$z>3$ auroral/\te{} anchor set beyond that number, then re-test the disputed claims
on matched-mass samples with \te{} anchors. ``Concrete, falsifiable, and the
measurement exists --- statistics are the gap.''

This paper measures the gap itself. Rather than assert that public statistics are
insufficient, we executed the assembly that a matched-mass re-test would require —
under a contract frozen before any science row was fetched — and we report what the
public archives actually supplied: how many galaxies at $z>3$ carry a
publicly tabulated auroral flux, how many of those survive a uniform, pre-registered
quality floor with the joins (H$\beta$, stellar mass) that a mass--metallicity
comparison needs, and how far the survivors fall below the pre-committed forecast.
The result is a census with an exact exclusion taxonomy, and a null stated in the
only form the contract licenses.

Two disciplines govern every sentence here. First, \emph{modality}: this is a census
and a null; no statement in this paper claims a metallicity deficit, a metallicity
evolution, or their absence. Second, \emph{provenance}: every number traces to a
sha-pinned lane artifact produced by a reviewed script; no number was typed by an
agent.

\section{The frozen measurement contract}\label{sec:contract}

The contract stack was written and frozen before the science fetch, in the order:
measurement design (adversarially refuted and patched), contract semantics, frozen
forecast, a single pre-result amendment ruling, and a frozen derivation pipeline.
We summarize the rules that shaped the census; the full semantics are in the lane
artifacts.

\subsection{Declared scale and decision rule}\label{sec:rule_s}

A metallicity value enters a comparison only on a \emph{declared scale}: named
method and calibration, an explicit reference frame or cited conversion, a
separately carried conversion-uncertainty term, and a channel label. Values failing
any condition are excluded by rule, not judgment. The decision rule (``Rule~S'')
bars the words ``detection'' or ``measured deficit'' for any offset that does not
exceed its own sample's scale uncertainty; for a pure direct-\te{} sample the
applicable scale floor is the 0.15 dex per-anchor \te{}-scale class. Scale terms are
frozen before the fetch and may not be recomputed smaller after a result exists.

\subsection{Anchor classes and the A$'$ amendment}\label{sec:aprime}

Sample members are classed by measurement content, never by the value obtained:
Class A (source-computed direct-\te{}), Class B (\te{}-consistent limits), Class C
(strong-line on a declared scale, never an anchor), Class X (excluded). During
catalog inspection — before any metallicity value existed in the lane — a conflict
surfaced: the JADES public tables publish auroral line \emph{fluxes} but no
source-computed \te{} or O/H, which the original Class A definition would discard.
A pre-result amendment ruling accepted \textbf{Class A$'$}: direct auroral
detection with lane-derived \te{}, under mechanical conditions: (i) eligibility
requires a source-published auroral flux and uncertainty giving
S/N$\geq5$ as tabulated, a floor declared pre-fetch against the documented
low-S/N $\lambda$4363 blending regime ([\ion{Fe}{2}] $\lambda$4360) and applied
uniformly; (ii) exactly one frozen derivation pipeline for all A$'$ members
(Sect.~\ref{sec:pipeline}); (iii) per-object provenance sufficient for independent
recomputation from the manifest row alone; (iv) A$'$ is a labelled sub-class with a
mandatory A-vs-A$'$ seam test; (v) equal standing with Class A in every rule.

The same ruling confirmed the lensing clause operating as written: a lensed sample
without per-object magnification $\mu$ and uncertainty defaults to
\emph{cluster-line-of-sight} and is excluded from matched-mass statistics — ``the
absence of magnification data is a declaration of exclusion, not an invitation to
assume $\mu\approx1$.'' The in-corpus $M_*\approx10^{5.7}\,M_\odot$ lensing-cluster
sample left the anchor set by this rule.

\subsection{The pre-committed forecast and the licensed null}\label{sec:forecast}

Before any science row was fetched, the expected \te{}-anchored yield was frozen
(forecast v1, \texttt{F4\_SHAPE2\_PREFETCH}, SHA-256
\texttt{61d48d22d34a2aed5fa1385a76afb04bb5aa6ac2f074c9b1b099961f99b860ec}): 10, 12,
and 3 anchors expected in the $8\leq\logm<9$, $9\leq\logm<10$, and $\logm\geq10$
bins respectively — \textbf{25 in total} — with a 0.25 dex deficit-precision
forecast and a 0.30 dex null-information threshold, based on catalog row counts and
known high-$z$ auroral detection rates ($\sim$1\% of deep spectroscopic samples).
Because the amendment changed the eligibility universe in both directions
(A$'$ added JADES-class candidates; the lensing rule removed the $10^{5.7}$ sample),
one re-freeze was licensed, still pre-fetch: forecast v2
(\texttt{F4\_SHAPE2\_PREFETCH\_V2}) expected 35/42/10 — \textbf{87 in total} — with
a 0.12 dex precision claim. Both versions are receipted; the contract requires any
null to cite the governing forecast \emph{and} disclose the supersession. The
adversarial reviewer's mock-run forensics had already flagged v2's 0.12 dex
precision claim as unachievable (finding F-T4-1); accordingly, the honest headline
of Sect.~\ref{sec:null} is stated against v1's \emph{more conservative} expectation.

The contract's null template (its \S6) is the only licensed form of a null result:
it must quote the frozen forecast, the realized statistics labelled as such, and an
exclusion bound — ``we saw nothing'' is barred, and so is any phrasing implying a
measurement of zero. A pre-committed minimum of 3 anchors per bin gates any per-bin
verdict.

\section{Census method}\label{sec:method}

\subsection{The reviewed-script protocol}\label{sec:protocol}

Every number in this paper was produced by a single script (\texttt{t3\_real.py}),
authored by one agent (Hwao) and line-by-line code-reviewed by an adversarial
reviewer (Kun) \emph{before} execution; no agent typed a number into any result
file. The protocol exists because an earlier executor seat in this lane produced a
stencilled mock result, which the reviewer's forensics rejected; the mock is
preserved for audit and contributed nothing to this paper. The review chain is
itself a reportable component of the method: an initial verdict of
\textsc{approved\_with\_edits} with three blocking edits (a S/N floor that
contradicted the frozen amendment; a doubly wrong dust-correction constant chain
that would have under-applied reddening corrections $\sim3\times$; an anchor-frame
table $\sim$0.13 dex low at the low-mass end), four required edits (ICF-fallback
exclusion, prediction-frame discipline, the 0.15 dex class floor, provenance
counts), and three advisories --- all applied pre-execution --- followed by six
reviewed micro-deltas during the live repair sequence (global enumeration;
sibling-redshift join; quote-layer repair; per-table fetch guards; a flux/error
column-pick repair; sibling-mass join). Seven runs are logged: three honest
empty-enumeration failures, one incomplete run, one completed run whose 748
candidate rows yielded zero derivable metallicities (traced by forensics to a
schema-order defect that had selected error columns as fluxes; repaired in the
fifth micro-delta), and two clean completions, the last adding only the mass join.

After the real run, the reviewer independently re-derived every published number:
all five anchors' $E(B-V)$, \te{}, and $\oh$ were reproduced \emph{to the printed
digit} from the archived fluxes with an independent PyNeb chain; all five
flux rows were verified character-for-character against the source archive; all
five stellar masses and redshifts were spot-verified live against the source's
sibling table; and the exclusion histogram was audited row-by-row (verdict:
\textsc{sound\_with\_corrections}; the two corrections repaired result-file
accounting and null framing, and changed no measured value).

\subsection{Enumeration and fetch}\label{sec:enumeration}

Enumeration was global and data-driven, not survey-assumed: a single TAP query over
the VizieR \texttt{TAP\_SCHEMA} found every table carrying a $\lambda$4363-class
column (79 tables), of which 23 also carry a $\lambda$5007-class column and a
redshift (in-table or via a sibling table join). At run time (2026-08-04, single
epoch), 12 of the 23 tables were unreachable (HTTP errors after four polite
retries each) and 11 tables in 8 catalogs were fetched. Rows were kept if the
auroral flux was tabulated and $z>3$: \textbf{95 rows}, from the ERO/GLASS/CEERS
auroral compilation (VizieR \dataset[J/ApJS/269/33]{https://vizier.cds.unistra.fr/viz-bin/VizieR?-source=J/ApJS/269/33},
10 rows) and the JADES public spectroscopic tables (VizieR
\dataset[V/159]{https://vizier.cds.unistra.fr/viz-bin/VizieR?-source=V/159}
\texttt{gngrat}/\texttt{gnprism}/\texttt{gsgrat}/\texttt{gsprism}; 9, 32, 6, and 38
rows). The reachable large low-$z$ surveys contributed zero $z>3$ rows, as
expected. Stellar masses were joined from the compilation's 182-row sibling table
(\texttt{tabled1}; 180 usable ID$\to\log M_*$ entries), with cross-table redshift
consistency verified.

\subsection{The frozen A$'$ derivation pipeline}\label{sec:pipeline}

One pipeline, frozen before the first \te{} was computed, applied identically to
every candidate: \te{}([\ion{O}{3}]) from the
$(\lambda4959+\lambda5007)/\lambda4363$ ratio via PyNeb \citep{Luridiana2015} at
$n_e=100\ \mathrm{cm^{-3}}$; the \citet{Izotov2006}-class $T(\mathrm{O^+}) =
0.7\,T_e + 3000$ K relation; ionic abundances from
$\lambda5007/\mathrm{H\beta}$ and $\lambda3727/\mathrm{H\beta}$ with
$\mathrm{ICF(O)}=1$ (O/H $=$ O$^+$/H$^+$ + O$^{2+}$/H$^+$); \citet{Cardelli1989}
reddening ($R_V=3.1$) from the source-published Balmer decrement
(H$\gamma$/H$\beta$, Case B intrinsic 0.468) with the
$\lambda4363$-vs-$\lambda5007$ differential correction taken from the same law
object; and a Monte Carlo uncertainty specification with fixed seed and draw count.
Rows lacking a tabulated auroral-flux uncertainty cannot demonstrate the floor and
are Class X by rule; rows whose $\mathrm{O^+}$ line is absent may not enter
compared statistics via a fabricated ionic fraction (none of the five anchors
needed either concession: all five have measured [\ion{O}{2}] $\lambda3727$, and
zero rows were ICF-fallback or missing-uncertainty). All five anchors' oxygen
abundances are referred, for display only, to the \citet{AndrewsMartini2013}
\te{}-anchored local mass--metallicity relation in its published asymptotic form
(their eq.~5) as the declared $z<3$ anchor frame; no strong-line value enters
anywhere in this paper.

\section{Results}\label{sec:results}

\subsection{The five contract-grade anchors}\label{sec:anchors}

Table~\ref{tab:anchors} lists every galaxy at $z>3$ for which, as of 2026-08-04,
the public archives supply a contract-grade direct-\te{} metallicity with a linked
stellar mass under a uniform pre-registered floor: five objects, all from the
ERO/GLASS/CEERS auroral compilation (three ERO, two GLASS), spanning
$z=4.015$--$8.496$, $\logm = 7.60$--$9.12$, and $\oh = 7.109$--$8.032$
(\te{} $= 14{,}307$--$24{,}847$ K). Figure~\ref{fig:gap}a shows them against the
declared anchor frame. Each row of Table~\ref{tab:anchors} was independently
recomputed from the archived fluxes by the adversarial reviewer and reproduced to
the printed digit; the five $E(B-V)$ values (0.181, 0.000, 0.046, one flagged
zero-reddening reading, 0.012) are varied and physical.

\begin{deluxetable*}{llccccccl}
\tablecaption{The five contract-grade public direct-\te{} anchors at $z>3$.\label{tab:anchors}}
\tablehead{
\colhead{ID} & \colhead{Source table} & \colhead{$z$} & \colhead{S/N($\lambda$4363)} &
\colhead{$E(B-V)$} & \colhead{$T_e$ (K)} & \colhead{$\oh$} & \colhead{$\logm$} &
\colhead{Notes}
}
\startdata
ERO\_04590   & J/ApJS/269/33/table1 & 8.496 & 5.11 & 0.181 & 24847.1 & 7.109 & 7.60 & below bin floor \\
ERO\_05144   & J/ApJS/269/33/table1 & 6.378 & 5.14 & 0.000 & 15456.7 & 7.922 & 8.55 & \\
ERO\_10612   & J/ApJS/269/33/table1 & 7.660 & 7.89 & 0.046 & 19391.9 & 7.685 & 7.78 & below bin floor \\
GLASS\_150029 & J/ApJS/269/33/table1 & 4.584 & 6.65 & \nodata & 16551.5 & 7.792 & 9.12 & dust corr.\ skipped\tablenotemark{a} \\
GLASS\_160133 & J/ApJS/269/33/table1 & 4.015 & 9.57 & 0.012 & 14307.4 & 8.032 & 8.11 & \\
\enddata
\tablenotetext{a}{Observed H$\gamma$/H$\beta$ $= 0.471 \geq 0.468$ (Case B): a
zero-reddening reading; no correction applied, flagged rather than silent.}
\tablecomments{All five have measured [\ion{O}{2}] $\lambda3727$ (ICF(O) $=1$ with
O$^+$ measured; no fallback rows). Masses from the 182-row sibling table
J/ApJS/269/33/tabled1, spot-verified live during forensics together with
cross-table redshift consistency. Every value was independently reproduced to the
printed digit from the archived fluxes.}
\end{deluxetable*}

\subsection{The exclusion census}\label{sec:exclusions}

Ninety of the 95 candidate rows were excluded, each with a per-row machine-recorded
reason (Table~\ref{tab:exclusions}); the full per-row file is a lane artifact. The
taxonomy is the census's substance --- it names exactly what the public archives
are missing:

\begin{deluxetable}{lrl}
\tablecaption{Exclusion taxonomy of the 95 public $z>3$ rows with a tabulated
$\lambda$4363 flux.\label{tab:exclusions}}
\tablehead{\colhead{Class} & \colhead{N} & \colhead{Dominant source}}
\startdata
Contract-grade anchors        & 5  & ERO/GLASS compilation \\
S/N($\lambda$4363) $<5$ floor & 64 & JADES prism + compilation \\
\quad of which zero tabulated flux & 14 & JADES prism \\
\quad of which near-miss (S/N 4.76--4.84) & 5 & mixed \\
No H$\beta$ $\to$ abundance undefined & 12 & JADES grating/prism \\
Missing $\lambda$4363/$\lambda$5007 flux & 6 & mixed \\
\te{} failure (non-physical ratio) & 8 & JADES grating/prism \\
\enddata
\tablecomments{Counts recomputed directly from the archived per-row sample file;
they sum to 95. Additional per-row flags: 20 rows carry no Balmer pair (no dust
correction possible, flagged); 1 anchor row has a flagged zero-reddening reading;
0 rows were excluded for missing tabulated flux uncertainty (Class X), 0 for
ICF fallback, and 0 anchors lacked a mass after the sibling join. The \te{}
failures are the physics working: e.g., one row has
$F(\lambda4363) > F(\lambda5007)$, an unphysical ratio for a genuine auroral
detection, and dies in the \te{} solver.}
\end{deluxetable}

Three features deserve emphasis. First, the S/N floor does the heaviest work (64
rows), and it is not a formality: the floor sits exactly in the documented
[\ion{Fe}{2}] $\lambda$4360 blending regime, and 14 of the excluded rows have a
tabulated flux of exactly zero. Second, the \textbf{near-miss band is real}: five
objects sit within 5\% of the floor (S/N $=4.76$--$4.84$), including ERO\_06355
($z=7.665$, $\logm=8.77$, S/N $=8.7/1.8=4.83$) and GLASS\_10021 ($z=7.286$,
$\logm=8.51$, S/N $=19.8/4.1=4.83$) --- both with complete Balmer sets and linked
masses, i.e., objects that would enter this census whole if their auroral
uncertainties improved by a few per cent. The floor bites at exactly 5.0 with no
fudge room; all 64 exclusions were recomputed directly from the archived
per-row fluxes, confirming every stated reason. Third, the JADES public tables ---
which contribute 85 of the 95 candidate rows and include auroral detections at
S/N up to 8.5 at $z=9.43$ --- contribute \emph{zero} anchors, not for want of
auroral photons but because the archived line tables carry no H$\beta$ columns
usable for abundances and no in-table stellar masses: twelve rows that pass the
S/N floor die at ``no H$\beta$ $\to$ abundance undefined,'' among them every
surviving $z>9$ candidate.

\subsection{Bins, the below-floor pair, and the null}\label{sec:null}

Applying the frozen bins: $N=2$ anchors in $8\leq\logm<9$ (ERO\_05144,
GLASS\_160133), $N=1$ in $9\leq\logm<10$ (GLASS\_150029), $N=0$ at $\logm\geq10$.
Two further verified anchors --- ERO\_04590 at $\logm=7.60$ and ERO\_10612 at
$7.78$ --- sit \emph{below the lowest frozen bin edge}: real, contract-grade, at
exactly the low masses this axis probes, and binned nowhere by the frozen 8.0
floor. They are counted here (and in the result file's
\texttt{below\_bin\_floor} field, added by forensic correction C1) so that no
reader concludes only three anchors exist. The accounting closes exactly:
$3$ binned $+\,2$ below-floor $=5$ anchors.

No bin reaches the pre-committed 3-anchor minimum, so the contract licenses
exactly one statement, quoted verbatim from the receipted result file (the \S6
template instantiation, with the required supersession disclosure of
Sect.~\ref{sec:forecast}):

\begin{quotation}
\noindent``Against the frozen forecast (v1 expected $\sim$25 usable \te{} anchors;
v2 expected $\sim$87), the executed public-archive assembly yielded 5
contract-grade anchors --- short of even the conservative expectation by roughly
an order of magnitude. No mass bin reaches the pre-committed 3-anchor minimum; no
deficit verdict is possible at contract-grade public statistics. This quantifies
A3's settle-line: the $z>3$ \te{} anchor set public archives can currently supply
is $N=5$.''
\end{quotation}

Figure~\ref{fig:gap}b displays the gap. Because the bins are unpopulated, no
per-bin offset, no exclusion bound at the 0.30 dex threshold, and no
scale-limited-vs-detection adjudication was computed or is quoted; the individual
anchor-minus-frame offsets visible in Figure~\ref{fig:gap}a are recorded in the
forensic ledger for provenance and carry no verdict.

\begin{figure*}
\centering
\includegraphics[width=\textwidth]{ANCHOR_GAP_FIGURE.pdf}
\caption{(a) The five contract-grade public direct-\te{} anchors at $z>3$
(Table~\ref{tab:anchors}) against the declared $z<3$ anchor frame
\citep[eq.~5 asymptotic form]{AndrewsMartini2013}; the shaded region marks masses
below the lowest frozen bin edge, where two of the five verified anchors fall.
No verdict is licensed by this panel: with every bin below the 3-anchor minimum,
a 0.15 dex per-anchor \te{}-scale floor, and an unresolved 0.14 dex anchor-frame
discrepancy (Sect.~\ref{sec:scope}), the visual offsets are not a measured
deficit. (b) The anchor gap: frozen forecast v1 (25 total), its licensed
pre-fetch revision v2 (87 total), and the realized contract-grade yield (2/1/0
per bin, plus 2 anchors below the lowest bin edge). No bin reaches the
pre-committed 3-anchor minimum (dotted line).\label{fig:gap}}
\end{figure*}

\subsection{What else the run refused to compute}\label{sec:refusals}

Two further honest refusals are part of the census. The FMR half of the parent
design (axis A4) is \texttt{not-computable-v1}: the fetched public tables declare
no SFR channel (an H$\beta$-based SFR would require a flux-calibration review that
the frozen pipeline does not license), and the script recorded the refusal rather
than improvise. And the confrontation of the eleven ledgered model/prediction
entries returned no verdicts: seven entries carry no numeric span in our declared
frame and four have no measured bins to confront; nothing defaults to
``consistent.''

\section{Discussion}\label{sec:discussion}

\subsection{What, specifically, would close the gap}\label{sec:close}

The exclusion census converts ``more data needed'' into named publication acts:

\begin{enumerate}
\item \textbf{Publish machine-readable flux tables with uncertainties and linked
masses.} The binding constraint is not auroral photons --- JADES alone contributes
85 public candidate rows, with auroral S/N up to 8.5 beyond $z=9$ --- but the
joins: archived line tables without usable H$\beta$ columns (12 kills above the
S/N floor, including every $z>9$ survivor) and without linked stellar masses.
A survey that archives its line fluxes \emph{with errors}, its H$\beta$, and an
ID-linked mass table converts its detections into contract-grade anchors at zero
telescope cost.
\item \textbf{The S/N $=4.8$ class.} Five objects sit within 5\% of the uniform
floor; two (ERO\_06355, GLASS\_10021) carry complete Balmer sets and masses and
would enter the census whole with marginally deeper integrations or improved
error models on exactly these lines. The near-miss band is where the anchor set
grows fastest per unit effort.
\item \textbf{Publish per-object magnifications.} The lensing rule removed the
$10^{5.7}\,M_\odot$ lensed sample not because its data are wrong but because the
public table carries no per-object $\mu$ with uncertainty. A citable companion
table of lens-model magnifications re-admits those objects under the contract's
existing rules.
\item \textbf{Archive availability.} Twelve of 23 candidate tables were
unreachable at run time (repeated HTTP errors). They are ``unreachable at run
time,'' not ``absent'': a retry pass can only add anchors, never remove the five
verified here.
\end{enumerate}

\subsection{Scope limits}\label{sec:scope}

This census is bounded by design, and the bounds are part of the result. The
enumeration is \emph{VizieR-only}: survey-team internal databases, value-added
reductions not deposited to VizieR, and papers whose measurements exist only in
manuscript tables are excluded \emph{by scope, not oversight} --- that is
precisely the point, since the settle-line requires data quotable at uniform
public rigor. The executed enumeration keyed on $\lambda$4363-class column names;
O~III] $\lambda$1666- and [\ion{N}{2}] $\lambda$5755-class anchors, though
A$'$-eligible in principle, were not enumerated in this first execution. The
census is a single epoch (2026-08-04). Finally, the declared anchor frame carries
a reportable internal discrepancy, quoted verbatim from the result file and
deliberately not averaged away: ``crew z9-10 receipted arithmetic implies
AM13(8.0)$\sim$8.26; this published form gives 8.12 ($\sim$0.14 dex) ---
REPORTABLE, referred to T4; not averaged away (Kun B3).'' With zero populated
bins this discrepancy is immaterial to every statement in this paper; it becomes
material the moment bins populate, and it is carried openly until adjudicated.

\subsection{Relation to the crew's $z\approx9$--$10$ unlensed study}\label{sec:z910}

This census is deliberately not a repeat of the crew's prior $z\approx9$--$10$
work, and the difference is the point. That study established the \emph{sign} of
a metal-poor offset with strictly unlensed field direct-\te{} anchors at
$z\approx9.3$--$10.6$ ($-0.69\pm0.03$ statistical against the same AM13 frame,
$N=5$ unlensed anchors, extended by GN-z11), and concluded --- with a total
systematic budget of $\pm$0.16 dex dominated by the same 0.15 dex \te{}-scale
class used here, bootstrap 95\% CI $[-0.82,-0.55]$ --- that the formal value is
\emph{systematic-limited, not a detection}. That discipline is the direct
precedent for this contract's Rule~S, and that study's lens-contamination repair
is the precedent for the lensing clause. What this census adds is orthogonal:
survey-wide public-archive \emph{scope} (a global enumeration rather than a
targeted anchor set) under a uniform contract-grade floor, quantifying what the
archives can supply across the full $z>3$ range. Notably, the executed census's
five anchors span $z=4.0$--$8.5$: no $z\gtrsim9$ candidate survived the public
joins (the $z=9.43$ pair falls at H$\beta$; the $z>9.6$ candidates lack the
nebular flux), so the two studies' anchor sets do not overlap --- targeted work
and archive-grade quotability are, today, disjoint. Both studies are honest about
the same absence: the FMR half was untouched there and is not-computable here;
the fixed-methodology FMR offset test at $z>3$ remains unexecuted by anyone,
including us.

\subsection{The consequence for the calibration-validity axis}\label{sec:a3}

Axis A3's binding statistic (c41\_012) counts $\sim$25 auroral-line galaxies at
$z>3$ in the corpus. This census measures the stricter quantity the settle-line
actually needs: how many of those detections are \emph{publicly quotable at
uniform rigor} --- tabulated fluxes with uncertainties, a uniform detection floor,
abundance-capable line sets, and linked masses. The answer is five, and zero
populated bins. The consequence is direct: the matched-mass re-test of c41\_043
against the c41\_033/035/044/045 cluster --- the named experiment that would
settle the axis --- cannot currently be run at contract grade from public
archives. Until the publication acts of Sect.~\ref{sec:close} occur, any
resolution of the calibration-validity dispute rests on data that are not
publicly quotable at uniform rigor, and referees and modelers alike should treat
apparent $z>3$ \te{}-anchor abundance samples of headline size accordingly: the
usable public core is an order of magnitude thinner than catalog row counts
suggest.

\section{Uncertainties and caveats}\label{sec:uncertainties}

\begin{enumerate}
\item \textbf{Archive reachability.} 12 of 23 candidate tables were unreachable
at run time; the census total (5) is therefore a floor with respect to fetch
failures --- a retry can only add anchors.
\item \textbf{Channel scope.} Only $\lambda$4363-class columns were enumerated;
A$'$-eligible UV-auroral ($\lambda$1666) and [\ion{N}{2}] $\lambda$5755 anchors,
if publicly tabulated anywhere with uncertainties, are not counted. Both this
and (1) bias the census \emph{conservative}: the true public yield can only be
$\geq5$, and the order-of-magnitude shortfall statement already absorbs this
direction.
\item \textbf{Anchor-frame discrepancy.} The 0.14 dex AM13 published-form vs
crew-arithmetic discrepancy (Sect.~\ref{sec:scope}) is carried unresolved. It
touches no statement here (no offsets are quoted) and shifts all per-object
display offsets coherently by $\sim$0.14 dex if resolved toward the crew
arithmetic.
\item \textbf{Per-object abundance uncertainties.} The frozen pipeline specifies
Monte Carlo propagation (fixed seed and draw count), but the v1 outputs tabulate
central values only; since no compared statistic was computed, no uncertainty
entered any verdict. Tabulated per-anchor uncertainties are required before any
populated-bin analysis, alongside the 0.15 dex per-anchor \te{}-scale class term.
\item \textbf{Pipeline version seam.} The pipeline freeze document names PyNeb
v1.1.18 atomic data; the executed and forensically reproduced environment carried
PyNeb 1.1.32. The independent reproduction to the printed digit was performed
under 1.1.32, so the seam is disclosed as a documentation defect, not a numeric
one. A second seam of the same class: the freeze document specifies the
H$\alpha$/H$\beta$ decrement as the primary dust input (source-published $A_V$
as fallback), whereas the executed and reviewed pipeline used the
H$\gamma$/H$\beta$ decrement exclusively --- at $z>4$ the only physically
available Balmer pair, so the primary branch was vacuous and the fallback never
fired; all five anchors were derived and reproduced under H$\gamma$/H$\beta$
consistently.
\item \textbf{Lensing metadata of the five anchors.} The anchors lie toward the
SMACS J0723 (ERO) and Abell 2744 (GLASS) sightlines, and the archived line table
publishes no per-object magnification column; the census consumes source-published
fluxes and masses as published, per the frozen A$'$ rule. Before these five
anchor any populated-bin statistic, a per-object magnification audit is required
--- the crew's own $z9$--$10$ lens-contamination repair is the standing precedent.
\item \textbf{One flagged zero-reddening reading.} GLASS\_150029's observed
H$\gamma$/H$\beta$ (0.471) sits marginally above the Case B intrinsic ratio; no
correction was applied and the row is flagged. Its abundance carries no dust bias
at the reported precision.
\end{enumerate}

\section{Summary}\label{sec:summary}

Under a measurement contract frozen before any science row was fetched --- a
declared \te{}-anchored scale, a uniform pre-registered S/N$(\lambda4363)\geq5$
floor, one frozen derivation pipeline, per-row exclusion provenance, and a
pre-committed yield forecast --- the public astronomical archives supply exactly
\textbf{five} contract-grade direct-\te{} metallicity anchors at $z>3$ (three ERO,
two GLASS; $z=4.015$--$8.496$), against a conservative frozen expectation of
$\sim$25 (and a licensed revised expectation of 87): a shortfall of roughly an
order of magnitude, with no stellar-mass bin reaching the pre-committed 3-anchor
minimum and two verified anchors falling below the lowest bin edge. Every anchor
was independently reproduced to the printed digit from the archived fluxes; all
90 exclusions carry audited per-row reasons. No deficit verdict of any size or
direction is possible at contract-grade public statistics, and the
calibration-validity dispute's resolution currently rests on data not publicly
quotable at uniform rigor. The census names the exact, cheap publication acts
that would change this: archive the fluxes with uncertainties, archive H$\beta$,
link the masses, publish per-object magnifications --- and re-observe or
re-reduce the five objects sitting within 5\% of the floor.

\begin{acknowledgments}
Pipeline authored by Hwao under the reviewed-script protocol; adversarial code
review and independent forensic reproduction by Kun; receipts and custody by
Tori; contract semantics and this draft by Lana. All artifacts, shas, and the
seven-run log are receipted in lane
\texttt{c41-trackb-shape2-mzr-20260804T1452K}. This draft is a lane artifact:
publication, database mutation, or promotion requires explicit human gating.
\end{acknowledgments}

\software{PyNeb \citep{Luridiana2015}, matplotlib, \texttt{nm\_external\_data}
(VizieR TAP client, cached and polite).}

\begin{thebibliography}{}

\bibitem[Andrews \& Martini(2013)]{AndrewsMartini2013} Andrews, B.~H., \&
Martini, P. 2013, ApJ, 765, 140

\bibitem[Cardelli et al.(1989)]{Cardelli1989} Cardelli, J.~A., Clayton, G.~C.,
\& Mathis, J.~S. 1989, ApJ, 345, 245

\bibitem[Izotov et al.(2006)]{Izotov2006} Izotov, Y.~I., Stasi\'nska, G.,
Meynet, G., Guseva, N.~G., \& Thuan, T.~X. 2006, A\&A, 448, 955

\bibitem[Luridiana et al.(2015)]{Luridiana2015} Luridiana, V., Morisset, C., \&
Shaw, R.~A. 2015, A\&A, 573, A42

\end{thebibliography}

% Source catalogs (cited in text by VizieR identifier; bibliographic entries to be
% completed at submission prep, outside this lane): J/ApJS/269/33 (ERO/GLASS/CEERS
% auroral compilation + tabled1 mass table), V/159 (JADES public spectroscopy).

\end{document}
